% [原文 p.19]
\begin{proof}[古典的证明]
由消去理论，齐次方程组存在公共解的条件，可写成与该方程组系数有关的若干个多项式（结式）等于 $0$（例如参看 [vdW91, 16.5]）。若多项式在无限多个素数 $p$ 处为 $0$，则它恒为 $0$。
\end{proof}

\begin{proof}[近代的证明]
齐次多项式组定义闭子概形 $Z\subset\PP^N_{\operatorname{Spec}R}$。投影 $\pi:\PP^N_{\operatorname{Spec}R}\to\operatorname{Spec}R$ 固有，故 $\pi(Z)\subset\operatorname{Spec}R$ 为闭子概形。若 $\pi(Z)$ 含有某个稠密子集的所有几何点，则 $\pi(Z)=\operatorname{Spec}R$。
\end{proof}

% [Tutor Note: 两个证明是同一件事的两种包装：结式（消去理论）把“有公共解”写成整系数行列式的消没，而 $\operatorname{Spec}R$ 上的闭子概形则把“几乎所有素数 $p$ 上可解”直接读成像集稠密。]

在现在的情形下，对几乎所有的 $p$，考虑关于变量 $(t_0,t_1)$ 的次数为 $m(\dim X+1)$ 的齐次形式 $(g_{p,0},\dots,g_{p,n})$ 中，伴随的有理映射 $\PP^1\dashrightarrow\PP^n$ 不是常值的那一些。其像 $C$ 含于 $X$ 中的条件，意味着
\[ h_i(g_{p,0},\dots,g_{p,n})=0\quad(\forall i) \]
是关于变量 $(t_0,t_1)$ 的恒等式。这一条件可看作 $g_k$ 的系数所满足的方程组。该方程组对属于某个 Zariski 稠密集的所有素数 $p$ 都有解，故由上述原理，它在任何代数闭域中都有解。

\begin{proof}[证明（续）]
\textbf{步骤 9}\quad 最后，注意步骤 2 与步骤 7 给出了过 $X$ 的任意点、次数不超过 $\dim X+1$ 的有理曲线。因此，若 $-K_X$ 丰富，便知 $X$ 被关于 $-K_X$ 次数至多 $\dim X+1$ 的有理曲线族所覆盖。
\end{proof}

\section{有理曲线的寻找（一般情形）}

\takeaway{本节只对某个固定的 $f$ 假设 $-(f(C)\cdot K_X)>0$。论证沿 §1.1 的框架：借正特征的 Frobenius 造出高次自同态，并在每次退化时记录曲线与 $H$ 的相交数，最终得到关于 $-K_X$ 次数和 $H$-次数都受控的有理曲线。}

\textbf{1.12}\quad 下面减弱关于 $X$ 的假设。即从今往后，不再假设 $-K_X$ 丰富，而只对所选定的某个 $f$ 假定
% [原文 p.20]
\[ -(f(C)\cdot K_X)>0. \]
另外，固定 $X$ 上的丰富除子 $H$。\hfill$\square$

接下来叙述的是主要结果。

\begin{theoremp}[{[Mor82]}]{1.13}
设 $X$ 为非奇异射影簇，$H$ 为 $X$ 上的丰富除子。若对某条不可约曲线 $C'\subset X$ 有 $-(C'\cdot K_X)>0$，则存在有理曲线 $E\subset X$ 满足
\[ \dim X+1\ge -(E\cdot K_X)>0,\qquad \frac{-(E\cdot K_X)}{(E\cdot H)}\ge\frac{-(C'\cdot K_X)}{(C'\cdot H)}. \]
\end{theoremp}

\begin{proof}
证明沿与前节相同的思路进行。不同之处在于，要对所得的曲线作更细致的考察。

\textbf{步骤 1}\quad 设 $f:C\to X$ 为 $C'$ 的正规化。若
\begin{equation}\tag{1.1}
-(f(C)\cdot K_X)-g(C)\cdot\dim X>0,
\end{equation}
则 $C$ 固定一点而形变。与以前一样，$f(C)$ 的形变必须退化为 $f'(C)+$（有理曲线之和）。

\textbf{步骤 2}\quad 为达到条件 (1.1)，移到正特征，并把 $f$ 与 Frobenius 态射的 $m$ 次幂复合。取 $m\gg0$，便可使 $p^m\cdot f(C_p)$ 退化为
\begin{equation}\tag{1.2}
(f^1_m)_*C_p+Z^1_{p,m}\equiv p^m\cdot f(C_p),
\end{equation}
其中 $Z^1_{p,m}$ 是有理曲线之和。定义正有理数 $M:=-(C'\cdot K_X)/(C'\cdot H)$。对几乎所有的 $p$，有
\[ M=-\frac{(f(C_p)\cdot K_{X_p})}{(f(C_p)\cdot H_p)}, \]
且该值在把 $f$ 换成 $f\circ F^m$ 后不变。若
\[ -((f^1_m)_*C_p\cdot K_{X_p})-g(C_p)\cdot\dim X_p>0, \]
则不必再把 $(f^1_m)_*C_p$ 与 Frobenius 态射复合，也能像以前一样使它移动。反复进行这些操作。每操作一次，$H_p$ 与 $(f^i_m)_*C_p$ 的相交数就减少，故该操作必在某个时刻停止。这样便得到数值等价
\[ (f^s_m)_*C_p+Z^s_{p,m}\equiv p^m\cdot f(C_p) \]
与不等式
% [原文 p.21]
\[ -((f^s_m)_*C_p\cdot K_{X_p})\le g(C_p)\cdot\dim X_p. \]
定义整数 $a,b,c,d$ 如下：
\begin{equation*}
\begin{aligned}
a&=-((f^s_m)_*C_p\cdot K_{X_p}), & b&=-(Z^s_{p,m}\cdot K_{X_p}),\\
c&=((f^s_m)_*C_p\cdot H_p), & d&=(Z^s_{p,m}\cdot H_p).
\end{aligned}
\end{equation*}
当 $m$ 很大时，$(c+d)$ 变大，而由 $(a+b)/(c+d)=M$ 知 $(a+b)$ 也变大。但 $a$ 有界，故 $b$ 变大。

% [Tutor Note: 与 §1.1 相比，这里多出的只有“记账”：每次退化都记录 $H$-次数，才能让正特征里的估计在回到特征 0 时仍给出与 $p$ 无关的界。]

\textbf{步骤 3}\quad 下面证明：对任意正数 $\epsilon$，当 $m$ 取得足够大时，$Z^s_{p,m}$ 的不可约成分 $E_p$ 中必有满足不等式
\begin{equation}\tag{1.3}
\frac{-(E_p\cdot K_{X_p})}{(E_p\cdot H_p)}>M-\epsilon
\end{equation}
者。

注意：若 $a/c<M$，则由引理 1.14 得 $b/d\ge M$，再由引理 1.14 便知所求的 $E_p$ 存在。又当 $m$ 取得很大时，若 $c$ 变大，则终有 $a/c<M$，由上述论证即可解决。若 $c$ 有界，则 $a$、$c$ 有界而 $b$、$d$ 变大，故 $b/d+\epsilon>(a+b)/(c+d)=M$ 成立。因此，取 $m\gg0$ 便有
\[ \frac{-(Z^s_{p,m}\cdot K_{X_p})}{(Z^s_{p,m}\cdot H_p)}>M-\epsilon. \]
于是由引理 1.14，所求的 $E_p$ 存在。

\begin{lemma}{1.14}
设 $c,d>0$。此时
\[ \frac{a+b}{c+d}\le\max\left\{\frac{a}{c},\frac{b}{d}\right\}. \]
\end{lemma}

% [Tutor Note: 引理 1.14 就是加权平均不等式：两个分式的加权平均不超过其中较大者。步骤 3 反复用它，保证退化成分中总有一个的斜率不低于 $M-\epsilon$。]

\textbf{步骤 4}\quad 假设 $-(E_p\cdot K_{X_p})>(\dim X+1)$。于是如定理 1.10 的步骤 2 所见，可固定两个以上的点而使有理曲线 $E_p$ 移动；这样曲线便退化为两条以上有理曲线之和。再次应用引理 1.14，可知退化成分 $E'_p$ 中至少有一个满足不等式 (1.3)。若再次
\[ -(E'_p\cdot K_{X_p})>(\dim X+1), \]
% [原文 p.22]
则 $E'_p$ 移动，如上所见可找到新的 $E''_p$ 满足不等式 (1.3)。然而每一步中 $(E_p\cdot H_p)$ 确实减小，故这一过程不能无限反复。最终可找到曲线 $E_p$（为简单起见仍记作 $E_p$）满足
\begin{equation}\tag{1.4}
\dim X+1\ge -(E_p\cdot K_{X_p})>0,\qquad \frac{-(E_p\cdot K_{X_p})}{(E_p\cdot H_p)}\ge M-\epsilon.
\end{equation}
这意味着 $(E_p\cdot H_p)\le(\dim X+1)/(M-\epsilon)$。式 (1.4) 中右式左端的分子与分母，都有与 $\epsilon$ 无关的上界。因此只要 $\epsilon$ 取得足够小，右式在令 $\epsilon=0$ 时不等式仍成立。

此外，$(E_p\cdot H_p)$ 有与 $p$ 无关的界。于是如定理 1.10 的步骤 8 那样，可证明复射影簇 $X$ 上存在满足定理 1.13 条件的有理曲线 $E$。
\end{proof}

\par\medskip\noindent\textbf{注意}\quad 这一论证无法对 $X$ 上有理曲线的位置说明任何事。不过，用别的论证可以知道：对 $C'$ 的任一点，都存在过它的有理曲线（参看 [MM86] 或 [Kol96, II.5]）。

\section{曲线锥（非奇异情形）}

\takeaway{本节引入 1-循环、$N_1(X)$ 与曲线锥 $\overline{NE}(X)$，并证明非奇异情形的锥定理：曲线锥由 $K_X\ge0$ 的部分与可数条极射线生成；把边界微移到 $K_X+\epsilon H\ge0$ 后，只剩有限多条极射线。证明的主干仍是 §1.2 的退化与计数。}

本节的主要目标，是证明非奇异射影簇的锥定理。这一结果由 [Mor82] 证明，是“森理论”主要的第一步。该定理的最终形式将在第 3 章确立。

\begin{definition}{1.15}
设 $K=\Q$ 或 $K=\R$，$V$ 为 $K$-向量空间。子集 $N\subset V$ 若满足 $0\in N$，且对正数倍封闭，则称为锥（cone）。

子锥 $M\subset N$ 若由 $u,v\in N$ 且 $u+v\in M$ 可推出 $u,v\in M$，则称为极端的（extremal）。$M$ 也称为 $N$ 的极锥面（extremal face）。一维极锥面也称为极射线（extremal ray）。

容易看出，不含过原点的直线的闭凸锥，是其所有极射线的凸包（convex hull）。
\end{definition}

% [原文 p.23]
\begin{definition}{1.16}
设 $X$ 为固有簇。1-循环（1-cycle）是不可约、既约的固有曲线的形式有限线性组合 $C=\sum a_iC_i$。系数全非负的 1-循环称为有效的（effective）。两个 1-循环 $C$、$C'$ 若对任何 Cartier 除子 $D$ 都有 $(C\cdot D)=(C'\cdot D)$，则称为数值等价的（numerically equivalent）。实系数 1-循环的数值等价类全体构成 $\R$-向量空间，记作 $N_1(X)$。1-循环 $C$ 所属的等价类记作 $[C]$。

（1）设 $NS(X)$ 为 $X$ 的 Néron–Severi 群（参看 [Har77]）。曲线 $C$ 与除子 $D$ 的相交数 $(C\cdot D)$（或 $(D\cdot C)$）给出非退化双线性形式
\[ N_1(X)\times\bigl(NS(X)\otimes_{\Z}\R\bigr)\to\R. \]
由 Néron–Severi 的底定理（theorem of the base of Néron–Severi），向量空间 $N_1(X)$ 是有限维的。证明请参看 [Kle66]。其维数称为 $X$ 的 Picard 数，记作 $\rho(X)$。

若 $X$ 在 $\C$ 上非奇异，则由关于 $(1,1)$-类的 Lefschetz 定理，存在单射 $N_1(X)\to H_2(X,\R)$（参看 [GH78, p.161]）。这说明：当 $X$ 在 $\C$ 上定义且非奇异时，$N_1(X)$ 是有限维的。

（2）设 $f:X\to Y$ 为固有态射。此时，1-循环 $Z$ 的顺像 $f_*Z$ 与 $f_*:N_1(X)\to N_1(Y)$ 如下定义。对任意不可约曲线 $C\subset X$，令
\[ \deg f|_C:=\begin{cases}\deg(f|_C:C\to f(C)), & f(C)\ \text{为曲线时},\\ 0, & f(C)\ \text{为点时},\end{cases} \]
并定义 $f_*C:=(\deg f|_C)\cdot f(C)$。对任意 1-循环 $Z=\sum a_iC_i$，由 $f_*Z:=\sum a_if_*C_i$ 扩张。注意：对 $Y$ 上的 Cartier 除子 $D$，由定义有 $(f_*Z\cdot D)=(Z\cdot f^*D)$。特别地，若 $Z\equiv0$ 则 $f_*Z\equiv0$，从而导出线性映射 $f_*:N_1(X)\to N_1(Y)$。还应注意：当 $X$、$Y$ 正规、$f$ 为双有理态射且 $Z$ 为 $\Q$-除子时，此 $f_*$ 与已有的定义一致。
\end{definition}

% [原文 p.24]
\begin{definition}{1.17}
设 $X$ 为固有簇。如下定义：
\begin{equation*}
\begin{aligned}
NE_{\Q}(X)&=\Bigl\{\textstyle\sum a_i[C_i]:C_i\subset X,\ 0\le a_i\in\Q\Bigr\}\subset N_1(X),\\
NE(X)&=\Bigl\{\textstyle\sum a_i[C_i]:C_i\subset X,\ 0\le a_i\in\R\Bigr\}\subset N_1(X),\\
\overline{NE}(X)&=NE(X)\ \text{在}\ N_1(X)\ \text{中的闭包}.
\end{aligned}
\end{equation*}
其中 $C_i$ 为 $X$ 上任意的不可约曲线。显然，$NE_{\Q}(X)$ 在 $\overline{NE}(X)$ 中稠密。

对任意除子 $D$，令 $D_{\ge0}:=\{x\in N_1(X):(x\cdot D)\ge0\}$（对 $>0$、$\le0$、$<0$ 也类似定义），并令 $D^{\perp}:=\{x:(x\cdot D)=0\}$。另外采用记号
\[ \overline{NE}(X)_{D\ge0}:=\overline{NE}(X)\cap D_{\ge0}, \]
对 $>0$、$\le0$、$<0$ 也用类似的记号。
\end{definition}

本书不展开一般理论，故关于态射与 $\overline{NE}(X)$ 的关系等，只在命题 3.27 中提及。考虑除子的锥同样自然，但本书不予处理。nef 除子锥的闭包与有效 1-循环锥互为对偶。该锥最初由 [Hir60] 处理。

下面叙述一般结果，并给出曲线锥的例子。最基本的结果来自 [Kle66]。证明在 §1.5 中给出。

\begin{theorem*}[{Kleiman 丰富性判别法（ampleness criterion）}]{1.18}
设 $X$ 为射影簇（或固有非奇异簇），$D$ 为任意 Cartier 除子。此时，$D$ 丰富的充要条件为
\[ D_{>0}\supset\overline{NE}(X)\setminus\{0\}. \]
\end{theorem*}

看图 1.1，就能很好地理解 $D$ 丰富时的状况。

\begin{figure}[htbp]
\centering
\resizebox{0.72\textwidth}{!}{\begin{tikzpicture}[line width=0.5pt]
\coordinate (O) at (0,0);
% 截口多边形（灰色区域）
\fill[gray!25] (-2.95,4.85) -- (0,5.55) -- (1.78,4.98) -- (2.58,4.30) -- (0.95,3.72) -- cycle;
% 从顶点 0 出发的射线
\draw (O) -- (-3.35,5.35);
\draw (O) -- (-2.62,5.30);
\draw (O) -- (0,7.15);
\draw[dashed] (O) -- (0.30,5.20);
\draw[dashed] (O) -- (1.10,4.30);
\draw (O) -- (2.10,5.95);
\draw (O) -- (2.95,4.95);
% 截口边界
\draw (-2.95,4.85) -- (0,5.55) -- (1.78,4.98) -- (2.58,4.30) -- (0.95,3.72) -- cycle;
\node[below left] at (O) {$0$};
% 曲线锥标号
\node (ne) at (2.45,6.55) {$\overline{NE}(X)$};
\draw (2.30,6.35) -- (1.95,5.45);
% 截口说明
\node[align=left,anchor=west] (cut) at (2.75,3.95) {超平面 $\{D=1\}$ 的\\截口是紧的};
\draw (2.70,4.08) -- (2.50,4.28);
\end{tikzpicture}
}
\figtitle{1.1}{曲线锥 $\overline{NE}(X)$ 与超平面 $\{D=1\}$ 的截口（原书 p.25）。}
\end{figure}

\begin{corollaryp}{1.19}
设 $X$ 为射影簇，$H$ 为丰富除子。此时下列成立。

（1）$\overline{NE}(X)$ 不含直线。

（2）对任何 $C>0$，集合 $\{z\in\overline{NE}(X):(z\cdot H)\le C\}$ 是紧的。

（3）对任何 $C>0$，整数系数的有效 1-循环 $Z=\sum a_iZ_i$ 中满足
% [原文 p.25]
$(Z\cdot H)\le C$ 者，只属于至多有限个数值等价类。
\end{corollaryp}

\begin{proof}
线性映射在去掉原点的直线上不可能取到真正为正的值。由此得到（1）。

其次，固定 $N_1(X)$ 上的范数 $\|\cdot\|$，假设 $W_C:=\{z\in\overline{NE}(X):(z\cdot H)\le C\}$ 不紧。则存在点列 $z_i\in W_C$ 使 $\|z_i\|\to\infty$。由于 $z_i/\|z_i\|$ 是范数为 1 的点列，适当取子列后，它收敛于某点 $y\in\overline{NE}(X)\setminus\{0\}$，且 $(y\cdot H)=\lim(z_i\cdot H)/\|z_i\|=0$ 成立。于是 $H$ 不丰富，矛盾。

最后，整数系数的 1-循环类在 $N_1(X)$ 中是离散的，故在任何紧集中至多只有有限个。
\end{proof}

% [Tutor Note: 推论 1.19 (2) 的紧性正是“$H$ 丰富”在曲线锥上的几何表现：超平面 $\{H=1\}$ 截出的集合有界，锥便不能在该方向上无限延伸。]

\begin{lemmap}{1.20}
设 $D$ 为固有不可约曲面上的除子，且 $(D^2)>0$。则当 $n\to\infty$ 时，$\dim|nD|\to\infty$ 或 $\dim|-nD|\to\infty$。
\end{lemmap}

\begin{proof}
由曲面情形的 Riemann–Roch 定理，有
\begin{equation*}
\begin{aligned}
h^0(nD)+h^0(K_X-nD)&\ge\frac{n^2}{2}(D^2)-\frac{n}{2}(D\cdot K_X)+\chi(\cO_X),\\
h^0(-nD)+h^0(K_X+nD)&\ge\frac{n^2}{2}(D^2)+\frac{n}{2}(D\cdot K_X)+\chi(\cO_X).
\end{aligned}
\end{equation*}
$n$ 取得很大时，上两式的右端都变大。但 $h^0(K_X-nD)$ 与 $h^0(K_X+nD)$ 不可能同时变大。这是因为两除子之和属于与 $n$ 无关的线性系 $|2K_X|$。故 $h^0(nD)\gg0$ 或 $h^0(-nD)\gg0$。此外，$n_1,n_2>0$ 时 $h^0(n_1D)>1$ 与 $h^0(-n_2D)>1$ 不会同时成立。（利用关系式 $n_2(n_1D)+n_1(-n_2D)=0$，可同上讨论。）故引理得证。
\end{proof}

% [原文 p.26]
\begin{corollaryp}{1.21}
设 $X$ 为不可约射影曲面，$H$ 为丰富除子。集合 $Q:=\{z\in N_1(X):(z^2)>0\}$ 有两个连通分支
\[ Q^+:=\{z\in Q:(z\cdot H)>0\},\qquad Q^-:=\{z\in Q:(z\cdot H)<0\}. \]
此外，$Q^+\subset\overline{NE}(X)$。
\end{corollaryp}

\begin{proof}
由 Hodge 指数定理，$N_1(X)$ 上作为二次形式的相交形式恰有一个正特征值。于是适当取基后，二次形式可写成 $x_1^2-\sum_{i\ge2}x_i^2$，并且可设 $H=(\sqrt{(H\cdot H)},0,\dots,0)$。由此得到 $Q$ 的两个连通分支
\[ Q^+=\Bigl(x_1>\bigl(\textstyle\sum_{i\ge2}x_i^2\bigr)^{1/2}\Bigr),\qquad Q^-=\Bigl(x_1<-\bigl(\textstyle\sum_{i\ge2}x_i^2\bigr)^{1/2}\Bigr). \]

对任何 $[D]\in Q$，$D$ 或 $-D$ 是有效的。而有效 1-循环与 $H$ 的相交数为正。因此 $Q$ 中的有效 1-循环类正是 $Q^+$ 的元素。
\end{proof}

% [Tutor Note: Hodge 指数定理把相交形式约化为只有一个正方向，故 $Q^+$ 就是以 $H$ 为轴的那一支；它恰好由有效曲线实现，从而与曲线锥的一支重合。]

\begin{lemmap}{1.22}
设 $X$ 为不可约射影曲面，$C\subset X$ 为不可约曲线。若 $(C^2)\le0$，则 $[C]$ 是 $\overline{NE}(X)$ 的边界元。若 $(C^2)<0$，则 $\R_{\ge0}[C]$ 是 $\overline{NE}(X)$ 的极射线。
\end{lemmap}

\begin{proof}
若 $D\subset X$ 是不可约曲线且 $(D\cdot C)<0$，则 $D=C$。若 $(C^2)=0$，则线性映射 $D\mapsto(D\cdot C)$ 在 $\overline{NE}(X)$ 上非负，而在 $C$ 上为 $0$。

这样，与 $C$ 不同的曲线的类都属于 $\overline{NE}(X)_{C\ge0}$，故 $\overline{NE}(X)$ 由 $\R_{\ge0}[C]$ 与 $\overline{NE}(X)_{C\ge0}$ 生成。若 $(C^2)<0$，则 $[C]\notin\overline{NE}(X)_{C\ge0}$，所以 $[C]$ 生成极射线。
\end{proof}

% [原文 p.27]
\begin{example}{1.23}
下面看几个例子。

（1）设 $X$ 为曲线 $B$ 上的极小直纹曲面。此时，$\overline{NE}(X)$ 是 $\R^2$ 的部分锥。于是由推论 1.19，$\overline{NE}(X)$ 恰由 2 条极射线生成。设 $f$ 为态射 $X\to B$ 的纤维的同调类，$s$ 为另一条极射线的生成元。由伴随公式（adjunction formula）得 $(f\cdot K_X)=-2<0$，再由推论 1.21 得 $s^2\le0$。

首先假设 $s^2<0$。此时可取收敛到 $\R_{\ge0}s$ 的有效 1-循环列 $D_n$。注意 $n\gg0$ 时 $(D_n^2)<0$。故 $\Supp D_n$ 的不可约成分 $E_n$ 中有满足 $(E_n^2)<0$ 者。于是由引理 1.22，$\R_{\ge0}[E_n]$ 是极射线；但除 $\R_{\ge0}f$ 之外的极射线只有 $\R_{\ge0}s$，故 $[E_n]\in\R_{\ge0}[s]$。

其次设 $s^2=0$。任取与纤维不同的不可约曲线 $D$。此时 $[D]$ 与 $f$ 生成 $N_1(X)$。考虑方程 $(xf+yD)^2=2xy(f\cdot D)+y^2(D\cdot D)=0$，则 $s$ 是 $2x(f\cdot D)+y(D\cdot D)=0$ 的解，其斜率为有理数。若 $B$ 的亏格为 0 或 1，则由直纹曲面的分类可知存在不可约曲线 $S\subset X$ 使 $[S]\in\R_{\ge0}[s]$。当 $g(B)\ge2$ 时，这样的曲线有时不存在。不过这种例子的构造并不简单。

例如，设表示 $\rho:\pi_1(B)\to U(2,\C)$ 使得 $\pi_1(B)$ 的任何有限指数子群的像都是 Zariski 稠密的。它对应 $B$ 上秩 2 的向量丛。此时取与之相伴的直纹曲面作为 $X$ 即可。详见 [Har70]。

% [Tutor Note: $s^2=0$ 时的斜率论证说明：$s$ 的方向由相交数唯一确定，必为有理方向；而能否由具体曲线实现，取决于基曲线 $B$ 的亏格。]

（2）设 $A$ 为 Abel 曲面，$H$ 为其上的丰富除子。Abel 曲面上任意曲线的自相交数非负，故用引理 1.20 可知 $\overline{NE}(A)=\overline{Q}^+$。若 $\dim N_1(A)\ge3$（例如 $A=E\times E$，$E$ 为某椭圆曲线），则 $\overline{NE}(A)$ 是“圆形”的锥。$\overline{NE}(A)$ 边界上的点都是极端的。这些点中大多数坐标不是有理数，因而不对应 $A$ 上的任何曲线。

（3）设 $X\subset\PP^3$ 为非奇异三次曲面。$\dim N_1(X)=7$，且 $X$ 上古典地已知有 27 条直线 $L_1,\dots,L_{27}$。由引理 1.22，它们生成 27 条极射线。

% [原文 p.28]
后文将看到
\[ NE(X)=\R_{\ge0}[L_1]+\cdots+\R_{\ge0}[L_{27}]. \]
特别地，$NE(X)=\overline{NE}(X)$ 是有限凸多面锥。对其他 Del Pezzo 曲面（即 $-K$ 丰富的曲面）也有同样的结论。

（4）（[Nag60]）取 $\PP^2$ 上三次曲线的束，使其每个成员都不可约。把它的 9 个基点爆破，所得曲面记为 $X$。由该束得到态射 $f:X\to\PP^1$，其一般纤维是椭圆曲线。从这 9 点中任选 1 点，把由它得到的 $f$ 的截面取为 0 截面，便在 $f$ 的纤维上引入群结构。若最初把束选得一般，则由其余 8 点得到的 8 个截面生成 $X$ 的无限自同构群。由此可知，$X$ 有无穷多条 $(-1)$-曲线，它们生成 $\overline{NE}(X)$ 的极射线。由于 $|-K_X|$ 是由椭圆曲线构成的束，$-K_X$ 是 nef 的，但不丰富。
\end{example}

把这些例子放在心上，就可以叙述森关于任意维数非奇异簇的最初结果了。有奇点的情形在第 3 章处理。非奇异情形的证明更几何化，值得先看一遍。

\begin{theorem*}[{锥定理（cone theorem），[Mor82]}]{1.24}
设 $X$ 为非奇异射影簇。

（1）存在至多可数个有理曲线 $C_i\subset X$，使 $0<-(C_i\cdot K_X)\le\dim X+1$，且
\[ \overline{NE}(X)=\overline{NE}(X)_{K_X\ge0}+\sum_i\R_{\ge0}[C_i]. \]

（2）对任何正数 $\epsilon>0$ 与丰富除子 $H$，下式成立：
\[ \overline{NE}(X)=\overline{NE}(X)_{(K_X+\epsilon H)\ge0}+\sum_{\text{有限和}}\R_{\ge0}[C_i]. \]
\end{theorem*}

% [Tutor Note: 锥定理 (2) 中“有限和”的意义在于：把边界从 $K_X\ge0$ 微移到 $K_X+\epsilon H\ge0$ 之后，剩下的负方向只剩有限多个，这是极射线局部有限性的雏形。]

这一定理的含义，看图 1.2 便可明了。

\begin{proof}
把数值等价的曲线视为相同，则满足 $0<-(C\cdot K_X)\le\dim X+1$ 的有理曲线 $C\subset X$ 至多可数。从每个等价类中选一条曲线 $C_i$。
% [边界: 定理 1.24 的证明起于 p.28、续到 ch01c；本块按块规则在此闭合，下一块以“证明（续）”接续。]
\end{proof}
